\chapter{Signals, computeds, effects}
\label{ch:signals}

The Elm runtime from Chapter~\ref{ch:elm} is one way to organise
state and updates. \maya{} ships a second, complementary mechanism:
\textbf{fine-grained reactivity} via \code{Signal<T>},
\code{Computed<T>}, and \code{Effect}. Inspired by SolidJS (and
Leptos in Rust), it lets you write code that \emph{looks} like
direct mutation but automatically tracks dependencies and runs only
the precise computations that depend on whatever changed.

This chapter explains the model, its data structures, and when to
reach for it over a \code{Program}.

\section{The problem reactivity solves}

In an Elm-style program, after every \code{Msg}, you call \code{view}
on the entire \code{Model} and re-render the whole UI. The diff
machinery makes this affordable, but the \emph{view function itself}
still touches every field --- even fields that did not change.

In a small app this is fine. In a 10,000-line dashboard with 200
panels, you would like a different property: when only one field
changes, only the views that read that field should re-evaluate.

That is what fine-grained reactivity gives you. It is not better
than Elm; it is a different point on the trade-off curve.

\section{The three primitives}

\subsection{\code{Signal<T>}: a value with subscribers}

A \code{Signal<T>} holds a value and remembers who is reading from
it:

\begin{lstlisting}
Signal<int> count{0};

int    n = count.get();        // read the current value
count.set(5);                  // write a new one
count.update([](int& x){ ++x; }); // mutate in place
\end{lstlisting}

Underneath, a \code{Signal} is a \code{ReactiveNode} (the base type
in \file{maya/include/maya/core/signal.hpp}) with a value field, a
list of \code{subscribers}, and a version counter.

\subsection{\code{Computed<T>}: a value derived from signals}

A \code{Computed<T>} is a function that uses other signals:

\begin{lstlisting}
Signal<int>   count{0};
Computed<int> doubled = computed([&] { return count.get() * 2; });

count.set(5);
int x = doubled.get();   // 10
\end{lstlisting}

Two ideas at play:

\begin{itemize}[leftmargin=*]
  \item When the lambda runs, every \code{get()} it calls is tracked
    as a \emph{dependency}. \maya{}'s reactive runtime keeps a
    thread-local ``current scope'' that signals push themselves
    onto.
  \item When any dependency changes, the computed is marked
    \emph{dirty}. The next \code{get()} on the computed re-runs the
    lambda; the get is cheap and idempotent.
\end{itemize}

\subsection{\code{Effect}: run a side effect when signals change}

An \code{Effect} is a lambda that runs once on construction and
re-runs whenever its dependencies change:

\begin{lstlisting}
Signal<int> count{0};
Effect e = effect([&] {
    std::println("count is now {}", count.get());
});

count.set(1);   // prints "count is now 1"
count.set(2);   // prints "count is now 2"
\end{lstlisting}

Effects are how reactive state reaches the outside world. The render
function in \maya{}'s \code{run(event\_fn, render\_fn)} quick-tool
API is effectively wrapped in an effect: the render runs once at
startup; subsequent renders happen automatically when any signal
read inside the render changes.

\section{How dependency tracking works}

Reactivity, when you understand it, is almost embarrassingly simple.

\begin{lstlisting}
namespace detail {
struct ReactiveNode {
    std::vector<ReactiveNode*> subscribers;
    std::vector<ReactiveNode*> dependencies;
    bool pending = false;       // O(1) de-dup flag
    uint64_t version = 0;
    virtual void mark_dirty() = 0;
    virtual void evaluate() = 0;
};

// Thread-local: the node currently being evaluated.
inline thread_local ReactiveNode* current = nullptr;
}
\end{lstlisting}

The dance:

\begin{enumerate}[leftmargin=*]
  \item When a computed or effect runs its lambda, it sets
    \code{current = this} for the duration.
  \item Inside the lambda, every \code{Signal<T>::get()} checks
    \code{current}. If it is non-null, the signal adds
    \code{current} to its \code{subscribers}, and the current node
    adds the signal to its \code{dependencies}.
  \item When \code{Signal<T>::set()} runs, it walks
    \code{subscribers} and calls \code{mark\_dirty()} on each.
  \item A dirty computed re-evaluates on next \code{get()}; a
    dirty effect schedules a re-run.
\end{enumerate}

The graph is bipartite-ish: signals are sources, effects are sinks,
computeds are both. Nothing is automatic but the bookkeeping; user
code just reads and writes signals as if they were variables.

\subsection{Batching}

If you write \code{count.set(5)} then \code{count.set(6)} in rapid
succession, you do not want effects to fire twice. \maya{} provides
\code{Batch}:

\begin{lstlisting}
{
    Batch batch;
    count.set(5);
    count.set(6);
    name.set("new");
}   // batch destructor flushes all dirty effects, exactly once each
\end{lstlisting}

Inside a batch, signal writes accumulate dirty flags on subscribers
but do not actually re-run them. When the batch is destroyed, every
dirty subscriber runs once. The \code{pending} bool on
\code{ReactiveNode} is the deduplication flag --- it ensures an
effect that depends on three batched signals fires once, not three
times.

\section{When to reach for reactivity}

A rule of thumb: reactivity is most useful when the UI has many
\textbf{independent} sub-views over a large state. A clock that
ticks every second alongside a chat scroll that updates on every
message alongside a status bar that watches a network connection
--- all three can be effects in their own right, each touching only
the signals they care about.

A pure \code{Program} is more useful when the state is small enough
to re-render in full every frame, when you want to record-and-replay,
or when you want to test the update function in isolation.

You can mix the two. \agentty{} runs Elm at the top level
(\code{Model} + \code{Msg} + \code{update}) and uses signals
internally for ephemeral UI state (cursor position, scroll offset,
animation phase).

\section{A worked example}

A weather panel that re-renders only the temperature line when the
temperature changes:

\begin{lstlisting}
Signal<double>    temperature{20.5};
Signal<int>       humidity{60};
Signal<std::string> location{"Tokyo"};

auto panel = v(
    t<"Weather"> | Bold,
    dyn([&] { return text("Location: "    + location.get()); }),
    dyn([&] { return text("Temperature: " + std::to_string(temperature.get())); }),
    dyn([&] { return text("Humidity: "    + std::to_string(humidity.get()) + "%"); })
);
\end{lstlisting}

The \code{dyn(...)} lambdas each read exactly one signal. Under
\maya{}'s reactive runtime, only the changed line is recomputed when
its signal changes. The diff machinery from Chapter~\ref{ch:simd}
then notices that only a few cells differ and emits a few bytes of
ANSI. The total cost of changing the temperature is on the order of
\emph{microseconds}.

\section{Lifetime and ownership}

A natural question: who owns the \code{ReactiveNode} graph? \maya{}'s
nodes are heap-allocated and reference-counted, but the references
are scoped:

\begin{itemize}[leftmargin=*]
  \item \code{Signal<T>} is a value handle. Copy it freely; the
    underlying node is shared.
  \item \code{Computed<T>} likewise.
  \item \code{Effect} is RAII. When the \code{Effect} object goes out
    of scope, the effect is unregistered from its dependencies and
    will not fire again.
\end{itemize}

The thread-local current scope means signals are \emph{not}
thread-safe across threads. The doc-level rule: one reactive graph
per thread. Cross-thread updates go through \code{Cmd<Msg>::task} (a
\code{Program} mechanism) or through a thread-safe queue you flush
on the UI thread.

\section{Comparison with React, Solid, Leptos}

If you have used reactive UI libraries:

\begin{itemize}[leftmargin=*]
  \item \maya{}'s model is closest to \textbf{Solid}'s. Signals are
    granular; tracking is automatic; there is no virtual DOM.
  \item Unlike \textbf{React}, there is no component re-render cycle.
    A signal change wakes only the effects that read it.
  \item Unlike \textbf{Leptos}, the host language is \cpp{} so
    there is no macro DSL --- signals are plain objects, the
    \code{|} pipes are the visual difference.
\end{itemize}

\begin{recap}
\maya{} ships a signal/computed/effect reactive runtime in
\file{core/signal.hpp}. Signals hold values; computeds derive;
effects run side effects. Dependency tracking is automatic via a
thread-local current-scope pointer. Batching deduplicates
notifications. Use reactivity for fine-grained sub-view updates;
use \code{Program} for testable top-level logic. The two compose.
\end{recap}

\section*{Workshop}

\begin{enumerate}[leftmargin=*]
  \item Write a clock effect that prints the current time every
    second. Drive it with a \code{Signal<std::string>} you update
    from a background thread (use \code{Cmd::Task} for the safe
    way).
  \item Build a \code{Computed<int>} that depends on two signals.
    Mutate them inside a \code{Batch} and verify the computed
    re-evaluates exactly once.
  \item Read \file{maya/include/maya/core/signal.hpp}. Find the
    \code{thread\_local} current-scope pointer. Trace how
    \code{Signal::get()} consults it.
\end{enumerate}
